% Generated by `taskgen course examples` from crypto.lfsr_recover (seed 3).
% Do not edit: change examples.yaml and regenerate.
\begin{example}[восстановление РСЛОС по $2L$ битам]\label{ex:lfsr-recover}
Известно, что последовательность $s_0, s_1, \dots$ над $\mathbb{F}_2$ порождается РСЛОС длины $L = 4$, т.\,е. $s_{n+4} = c_1 s_{n+3} + \dots + c_{4} s_n$. По первым $2L = 8$ битам \[1,\ 0,\ 0,\ 0,\ 1,\ 0,\ 1,\ 0\] восстановить коэффициенты $c_1, \dots, c_{4}$ (составить и решить линейную систему над $\mathbb{F}_2$ методом Гаусса) и найти следующие 4 бита.

\par\noindent\textit{Решение.} Соотношения при $n = 0, \dots, 3$ дают систему из $4$ уравнений: $c_{4} s_n + c_{3} s_{n+1} + \dots + c_1 s_{n+3} = s_{n+4}$. Строки расширенной матрицы (столбцы $c_{4}, \dots, c_1$ \textbar\ правая часть) --- окна длины $5$ последовательности:\[\left(\begin{array}{cccc|c}1 & 0 & 0 & 0 & 1 \\ 0 & 0 & 0 & 1 & 0 \\ 0 & 0 & 1 & 0 & 1 \\ 0 & 1 & 0 & 1 & 0\end{array}\right)\]Метод Гаусса над $\mathbb{F}_2$ (сложение строк = XOR):\par\noindent столбец $c_{4}$: преобразований не нужно;\par\noindent столбец $c_{3}$: $R_{2} \leftrightarrow R_{4}$ $\Rightarrow$ $\left(\begin{array}{@{}cccc|c@{}}1 & 0 & 0 & 0 & 1 \\ 0 & 1 & 0 & 1 & 0 \\ 0 & 0 & 1 & 0 & 1 \\ 0 & 0 & 0 & 1 & 0\end{array}\right)$;\par\noindent столбец $c_{2}$: преобразований не нужно;\par\noindent столбец $c_{1}$: $R_{2} \mathrel{+}= R_{4}$ $\Rightarrow$ $\left(\begin{array}{@{}cccc|c@{}}1 & 0 & 0 & 0 & 1 \\ 0 & 1 & 0 & 0 & 0 \\ 0 & 0 & 1 & 0 & 1 \\ 0 & 0 & 0 & 1 & 0\end{array}\right)$;\par Ранг равен $4$, решение единственно: $(c_1, \dots, c_{4}) = (0, 1, 0, 1)$, $f(x) = x^{4} + x^{2} + 1$. Проверка: рекуррентность порождает данные $8$ бит; продолжение: \texttt{1000 1010 0010} (группы по $L$ бит).\[f(x) = x^{4} + x^{2} + 1,\quad s_{8}, \dots, s_{11} = 0,\ 0,\ 1,\ 0.\]
\end{example}
